\section{Method}
\label{sec:method}

We begin by formalizing the frozen-feature classification pipeline, then describe each dimensionality reduction method evaluated in this study, and finally present two extensions we propose.

\subsection{Problem Setting}

Let $f_\theta: \mathcal{X} \to \mathbb{R}^D$ denote a CNN backbone with frozen weights $\theta$, pretrained on ImageNet~\cite{deng2009imagenet}. Given a target dataset $\{(\mathbf{x}_i, y_i)\}_{i=1}^{N}$ with $C$ classes, we extract features $\mathbf{z}_i = f_\theta(\mathbf{x}_i) \in \mathbb{R}^D$ and seek a projection $\mathbf{W} \in \mathbb{R}^{D \times d}$ (with $d \ll D$) such that a linear classifier trained on the projected features $\mathbf{W}^\top\mathbf{z}_i \in \mathbb{R}^d$ achieves high accuracy on held-out data.

The pipeline is: \emph{Image} $\xrightarrow{f_\theta}$ \emph{Feature} ($D$-dim) $\xrightarrow{\mathbf{W}}$ \emph{Reduced Feature} ($d$-dim) $\xrightarrow{g}$ \emph{Prediction}.

For the classifier $g$, we use $\ell_2$-regularized logistic regression throughout, with identical hyperparameters across all methods. This isolates the effect of the dimensionality reduction step.

\subsection{Methods Compared}

We compare ten methods organized into four categories. All methods except \emph{Full} produce $d = C - 1$ dimensional features (99 for CIFAR-100, 199 for Tiny ImageNet), unless otherwise noted.

\subsubsection{Control Methods}

\paragraph{Full Features} The feature vector $\mathbf{z}_i$ is passed directly to the classifier without reduction. This is the standard transfer learning baseline.

\paragraph{PCA} Principal Component Analysis~\cite{pearson1901pca} projects onto the $d$ leading eigenvectors of the feature covariance matrix. PCA is unsupervised and serves as a lower bound for what label-free dimensionality reduction can achieve.

\subsubsection{Classical Supervised Methods}

\paragraph{LDA} Linear Discriminant Analysis~\cite{fisher1936lda} solves the generalized eigenvalue problem $\mathbf{S}_b\mathbf{w} = \lambda\mathbf{S}_w\mathbf{w}$, where
\begin{equation}
\mathbf{S}_b = \sum_{c=1}^{C} n_c(\boldsymbol{\mu}_c - \boldsymbol{\mu})(\boldsymbol{\mu}_c - \boldsymbol{\mu})^\top, \quad
\mathbf{S}_w = \sum_{c=1}^{C}\sum_{i \in \mathcal{C}_c} (\mathbf{z}_i - \boldsymbol{\mu}_c)(\mathbf{z}_i - \boldsymbol{\mu}_c)^\top,
\end{equation}
with $\boldsymbol{\mu}_c$ and $\boldsymbol{\mu}$ the class and global means, $n_c = |\mathcal{C}_c|$ the class count, and $\mathcal{C}_c$ the set of indices belonging to class $c$. The $d = C - 1$ eigenvectors corresponding to the largest eigenvalues form the projection $\mathbf{W}$. We use the implementation in scikit-learn~\cite{sklearn}, which applies an SVD-based solver that handles rank-deficient scatter matrices gracefully.

\paragraph{PCA+LDA} A two-stage approach: PCA first reduces dimensionality to $C - 1$, then LDA is applied in the reduced space. This is sometimes recommended when $D \gg N$ to improve the conditioning of the scatter matrices.

\subsubsection{Academic Variants}

\paragraph{Regularized LDA (R-LDA)} Replaces $\mathbf{S}_w$ with the shrinkage estimate~\cite{friedman1989regularized}:
\begin{equation}
\hat{\mathbf{S}}_w = (1 - \alpha)\mathbf{S}_w + \alpha \cdot \frac{\tr(\mathbf{S}_w)}{D}\mathbf{I},
\end{equation}
where the shrinkage intensity $\alpha$ is set automatically using the Ledoit--Wolf estimator~\cite{ledoit2004well}. This stabilizes inversion when $\mathbf{S}_w$ is near-singular, which can occur with high-dimensional features and limited samples per class.

\paragraph{Local Fisher Discriminant Analysis (LFDA)} Sugiyama's extension~\cite{sugiyama2007lfda} constructs localized scatter matrices using a $k$-nearest-neighbor affinity matrix, preserving multimodal class structure. We set $k = 7$ following the original recommendation.

\paragraph{Neighbourhood Components Analysis (NCA)} Goldberger et al.'s metric learning approach~\cite{goldberger2004nca} learns a projection by maximizing a stochastic $k$-NN leave-one-out accuracy objective. We use the scikit-learn implementation with default hyperparameters.

\subsubsection{Proposed Extensions}

We additionally evaluate two lightweight extensions of LDA that exploit complementary information in the feature space.

\paragraph{Residual Discriminant Augmentation (RDA)} LDA projects onto $C-1$ directions that maximize class separation. The orthogonal complement---the space \emph{not} captured by LDA---may still contain useful variance. RDA appends a small number of PCA components computed in this residual subspace:
\begin{equation}
\mathbf{z}_i^{\text{RDA}} = \left[\mathbf{W}_{\text{LDA}}^\top\mathbf{z}_i \ ;\ \mathbf{W}_{\text{res}}^\top(\mathbf{z}_i - \mathbf{W}_{\text{LDA}}\mathbf{W}_{\text{LDA}}^\top\mathbf{z}_i)\right] \in \mathbb{R}^{(C-1)+k},
\end{equation}
where $\mathbf{W}_{\text{res}}$ contains the $k$ leading PCA components of the LDA residuals, with $k$ set to 20 or 30 depending on the original feature dimension. The total dimensionality is $(C-1) + k$, still far below $D$.

\paragraph{Discriminant Subspace Boosting (DSB)} Inspired by boosting~\cite{freund1997decision}, DSB iteratively reweights the training samples and refits LDA. In each round, samples that were misclassified in the previous round receive higher weight in the scatter matrix computation. After $T$ rounds, the final projection is the LDA solution from the last round. We use $T=2$ rounds, which balances the small accuracy gain against the doubling of LDA fitting time. The projection dimensionality remains $C-1$.

Both extensions are computationally inexpensive relative to methods like NCA or LFDA, adding only a constant factor to LDA's fitting time.

\subsection{Standardization}

All features are standardized to zero mean and unit variance before classification, using statistics computed on the training set. This is critical for fair comparison: different reduction methods produce features at different scales, and the $\ell_2$-regularized logistic regression is sensitive to feature scaling. Standardization is applied \emph{after} projection, following standard machine learning practice.
